\documentstyle[12pt,fullpage,steve]{article}
\setcounter{page}{1}

\begin{document}

\title{EE 150 Quiz \#2}
\author{Due: Oct. 16, 1996 2:40 pm (in class)}
\date{}
\maketitle

In a popular board game, three of the characters include Colonel
Mustard, Mrs. White, and Professor Plum.  Imagine you have to write a
function {\tt return-the-other-one} which accepts as parameters {\bf two}
{\tt chars}, each of which is a 'm' (Mustard), 'w' (White) or 'p'
(Plum).  This function returns the letter corresponding to the
character {\em not} passed as a parameter.  Note that the characters
are passed in no particular order.

For example, a call to the function {\bf return-the-other-one( 'm',
'w')} results in a returned value of 'p'.  

Given the following skeleton for the function, give an implementation
using ANDs (\&\&) and ORs ($||$) of each of the conditions 1 and 2
that will make the following function behave as advertised.  You may
not change the structure of the function, only specify the appropriate
conditions. 
\begin{verbatim}
char return-the-other-one( char c1, char c2 )
{
if ( ...condition 1... )
return 'p';               /* then the other character is Professor Plum */
else
if ( ...condition 2... )
return 'w';               /* then the other character is Mrs. White */
else
return 'm';               /* last possibility is Colonel Mustard */  
}
\end{verbatim}

\end{document}
